
\emph{suppose que pour tout point @@M0@@, et toute composante irr\'eductible @@M1@@ de @@M2@@, les composantes irr\'eductibles de @@M3@@ soient de dimension @@M4@@, (ce qui est le cas si les composantes irr\'eductibles des fibres de @@M5@@ sont de dimension @@M6@@). Alors le morphisme induit entre les foncteurs de Picard}
@@M7@@
\emph{est de type fini.}

\emph{De plus, si @@M8@@ est int\`egre, il existe un ouvert non-vide @@M9@@ de @@M10@@ tel que les foncteurs de Picard restreints \`a @@M11@@ soient repr\'esentables et le morphisme induit}
@@M12@@
\emph{soit quasi-affine (de type fini).}

D\'emontrons d'abord la premi\`ere assertion. Comme dans la preuve de 3.6, on voit par r\'ecurrence noeth\'erienne qu'on peut remplacer @@M13@@ par un ouvert non-vide quelconque, et on se ram\`ene facilement au cas o\`u @@M14@@ est int\`egre et @@M15@@ est int\`egre sur @@M16@@, toujours en rempla\c{c}ant @@M17@@ par @@M18@@ lorsqu'on remplace @@M19@@ par @@M20@@\,; ceci fait, on se ram\`ene de la m\^eme fa\c{c}on au cas o\`u @@M21@@ est normal sur @@M22@@, en utilisant \`a la place de XII 1.3 le lemme suivant.

\textbf{Lemme 3.9.} \emph{Soient @@M23@@ un sch\'ema noeth\'erien int\`egre et @@M24@@ un @@M25@@-sch\'ema, int\`egre sur @@M26@@. Alors il existe un ouvert non-vide @@M27@@ de @@M28@@, un sch\'ema int\`egre @@M29@@ fini et fid\`element plat sur @@M30@@, un @@M31@@-sch\'ema @@M32@@ normal et int\`egre sur @@M33@@, et un @@M34@@-morphisme fini surjectif @@M35@@.}

